\chapter{Use and Functionality of the System}

This chapter provides a brief description of how the system is started and of the
functionalities it implements.


\section{Starting the System}

In order to run the expert system it is sufficient to copy the \texttt{PESAD}
folder, which contains the modules that implement the system, into a directory of
choice. The system relies on the SWI-Prolog environment and is launched by
consulting the start-up file \texttt{start.pl}, located in the \texttt{PESAD}
folder, directly from the command line:

\begin{lstlisting}
$ swipl start.pl
\end{lstlisting}

The very first action performed at start-up is the selection of the interface
language. The system asks the user to choose between English and Italian, and the
whole session is then conducted in the selected language:

\begin{lstlisting}
Select language / Seleziona lingua:
	1: English
	2: Italiano
 > 1.
\end{lstlisting}

Once the language has been chosen, the system prints a short presentation message
followed by the list of available commands, and then waits for the user to type a
command:

\begin{lstlisting}
-------------------------------------------------------------------
-- PESAD - Prolog Expert System for Anxiety,
--         Obsessive-Compulsive and Trauma-Related Disorders
--
--	knowledge base: DSM-5-TR
--	author:         Donato Meoli
-------------------------------------------------------------------

Available commands:

	investigation.	- Run the full investigation of all disorders;
	control.	- Run the check of a single disorder;
	facts.		- Show the asserted facts;
	help.		- Show the list of available commands;
	clean.		- Clear the memory;
	quit.		- Close the interface.

>
\end{lstlisting}

At this point the system remains idle, waiting for a user command. Once a command
has been typed, the system performs a short validity check on it: if the command
is not recognised, an error message is shown and the system waits again for a
valid command. As a general rule, every input typed by the user is followed by
rigorous validity checks.


\section{System Commands}

The expert system provides six different commands:

\begin{itemize}
  \item \texttt{investigation} --- starts the complete investigation of all the
        disorders represented in the system;
  \item \texttt{control} --- starts the check of a single disorder, chosen among
        those implemented in the system;
  \item \texttt{facts} --- displays the facts temporarily asserted by the user
        during the current consultation;
  \item \texttt{help} --- displays the list of commands accepted by the system;
  \item \texttt{clean} --- clears the working memory;
  \item \texttt{quit} --- closes the user interface.
\end{itemize}

The command words are always in English, regardless of the selected interface
language: only the questions, diagnoses and explanations that the user reads are
localised, whereas the commands are control tokens and are kept uniform.

Throughout a consultation, when the inference engine lacks the information needed
to continue its reasoning, it poses questions to the user. Symptom questions are
answered with \texttt{yes} or \texttt{no}; in addition, the user may type
\texttt{why} (or \texttt{perche} in the Italian interface) to obtain an
explanation of why the question is being asked. After a \texttt{yes}/\texttt{no}
answer, the system asks for a certainty value, which must be a real number in the
interval $[0.0, 1.0]$ expressing how confident the user is in the given answer.


\subsection{Investigation}

The \texttt{investigation} command launches the procedure that investigates the
symptoms affecting a given patient, in order to diagnose a list of possible
anxiety disorders. The result of the investigation is therefore an \emph{ordered}
list of disorders that are candidates to be the cause of the patient's distress.

This functionality is characterised by two design choices. The first is to
structure the reasoning so that, once an Axis~I disorder has been resolved, the
system does not stop the exploration but continues to examine alternative paths.
In this way the inference engine produces a list of anxiety disorders, each
associated with a certainty score. This score, obtained as the product of the
certainty of the premise and the certainty of the rule representing the disorder,
allows a \emph{ranking} of the anxiety disorders, from the most certain to the
least certain.

The second choice is to report in the result list all the disorders that were
examined, together with their certainty degree. This was considered appropriate
because the expert user of the system can benefit from a complete view of the
anxiety disorders and, whenever a discrepancy with their own reasoning is
observed, can investigate the reason for the difference by displaying the proof
tree of the disorders involved.

In the code, this functionality is implemented through the following predicate:

\begin{lstlisting}[language=Prolog]
solve(investigation) :-
    delete_trees,
    investigate_goals(Goals_List),
    descending_order_goals(Goals_List, Sorted_Goals_List),
    display_investigation(Sorted_Goals_List),
    investigation_how(Sorted_Goals_List).
\end{lstlisting}

The predicate \texttt{delete\_trees} removes any trace of the chain of rules and
facts that led to the various goals in previous executions. This operation does
\emph{not} remove the facts dynamically asserted by the user during previous
consultations: the two operations are deliberately kept separate, so that all the
facts stated by the user remain in memory across subsequent executions, while the
traces of the rules and facts used in the proof trees are removed at every run.
This choice prevents partial traces of disorders from remaining in memory should
the investigation procedure terminate prematurely because of unexpected errors or
incorrect use of the system. If the user wishes to clear the entire working
memory, the \texttt{clean} command can be used.

The predicate \texttt{investigate\_goals/1} first determines the patient's age
class --- child, adolescent or adult --- and then, for every disorder
represented, verifies the required features, asking the user questions whenever
the available information is insufficient. The predicate
\texttt{descending\_order\_goals/2} sorts the resolved goals in descending order
of certainty, \texttt{display\_investigation/1} prints the resulting list of
disorders together with their certainty degree, and
\texttt{investigation\_how/1} manages the optional request for an explanation of
how a conclusion was reached. A single derivation tree is printed at a time, to
make it easier to read: after the diagnosis has been shown, the system asks the
user whether they want the ``how''; if the answer is \texttt{yes}, the system
asks for the code of the disorder for which the explanation should be produced,
prints the corresponding tree, and then allows the user either to leave the
investigation procedure or to request further explanations for other disorders.

A concise example of an investigation session is shown below:

\begin{lstlisting}
> investigation.

[ ... the system asks the relevant symptom questions ... ]

Does the patient experience recurrent and unexpected panic attacks?
? (yes, no) > yes.
Enter the certainty degree (from 0.0 to 1.0) > 0.9.

Does the patient avoid situations from which escape might be difficult?
? (yes, no) > why.
Pervasive fear or avoidance of situations from which escape might be
difficult would direct the diagnosis in favor of agoraphobia.
? (yes, no) > yes.
Enter the certainty degree (from 0.0 to 1.0) > 0.8.

========================== DIAGNOSIS ==========================

1) Panic Disorder with a certainty degree of 81%

2) Agoraphobia with a certainty degree of 64%

Do you want to know how they were deduced? (yes, no) > yes.

Enter the code corresponding to the chosen option > 1.

[ ... the proof tree for Panic Disorder is printed ... ]

Do you want to know how they were deduced? (yes, no) > no.
\end{lstlisting}


\subsection{Control}

The \texttt{control} command launches the procedure that checks whether the
patient suffers from a given anxiety disorder. This functionality is of
fundamental importance whenever the expert has already carried out a preliminary
analysis of the patient, or has data that allow a particular type of disorder to
be hypothesised: with this tool the expert can confirm or reject that hypothesis.

In the code, the control of a given disorder is implemented through the following
predicate:

\begin{lstlisting}[language=Prolog]
solve(control) :-
    selection_goal(Goal_Name),
    delete_tree(Goal_Name),
    solve_goal(Goal_Name, CF),
    display_control(Goal_Name, CF),
    control_how(Goal_Name).
\end{lstlisting}

The predicate \texttt{selection\_goal/1} lets the user select the disorder to be
checked from a list of possible options, retrieved from the up-to-date list of
disorders present in the knowledge base. The predicate
\texttt{delete\_tree/1} then removes any trace of the chain of rules and facts
that led to the conclusion of that goal in a previous execution; as in the
investigation case, this operation removes only the proof-tree traces and not the
facts asserted by the user, again to avoid leaving partial traces in memory if
the procedure terminates prematurely. The entire working memory can still be
cleared with the \texttt{clean} command.

The predicate \texttt{solve\_goal/2} is very similar to
\texttt{investigate\_goals/1}, the difference being that the inference engine
verifies exclusively the disorder unifying with \texttt{Goal\_Name}; as before,
the reasoning is driven by the patient's age class, and questions are posed to the
user only when new information is needed. The predicate returns the certainty
degree with which the disorder is found to be true. The predicate
\texttt{display\_control/2} prints the single result of the check, and
\texttt{control\_how/1} asks the expert whether they are interested in an
explanation of how the conclusion was reached: if the answer is \texttt{yes}, the
system immediately prints the corresponding proof tree; otherwise it closes the
control procedure and waits for new commands.

A concise example of a control session is shown below:

\begin{lstlisting}
> control.

Which of the following disorders do you want to check?

	1: Panic Disorder
	2: Agoraphobia
	3: Generalized Anxiety Disorder
	[ ... ]

Enter the code corresponding to the chosen option > 1.

[ ... the system asks only the questions relevant to Panic Disorder ... ]

Does the patient experience recurrent and unexpected panic attacks?
? (yes, no) > yes.
Enter the certainty degree (from 0.0 to 1.0) > 0.9.

========================== DIAGNOSIS ==========================

The patient suffers from Panic Disorder with a certainty degree of 81%

Do you want to know how it was deduced? (yes, no) > yes.

[ ... the proof tree for Panic Disorder is printed ... ]
\end{lstlisting}


\subsection{Facts}

The \texttt{facts} command displays the facts that the user has temporarily
asserted through the interface during the current consultation. Each fact is
printed in the form \texttt{(Answer, Attribute, Value, Certainty)}, where:

\begin{itemize}
  \item \texttt{Answer} is the type of response, \texttt{yes} or \texttt{no}. For
        facts asserted through binary (yes/no) questions, the value is
        \texttt{yes} when the patient answered affirmatively and \texttt{no} when
        the answer was negative. For facts asserted through multiple-choice
        questions the value is always \texttt{yes}, since the answer corresponds
        to the option selected by the user;
  \item \texttt{Attribute} is the symptom or feature being investigated;
  \item \texttt{Value} is the item or option associated with the answer;
  \item \texttt{Certainty} is the certainty value, in $[0.0, 1.0]$, that the user
        associated with the answer.
\end{itemize}

An example of the output produced by this command is the following:

\begin{lstlisting}
> facts.

Asserted facts:

(yes, panic_attacks, recurrent_unexpected, 0.9)
(no, medical_condition, present, 0.95)
(yes, avoidance, social_situations, 0.8)
\end{lstlisting}


\subsection{Help}

The \texttt{help} command reprints the menu listing all the commands accepted by
the system, together with a short description of each, exactly as shown at
start-up. It is useful as a quick reminder of the available functionalities
during a consultation:

\begin{lstlisting}
> help.

Available commands:

	investigation.	- Run the full investigation of all disorders;
	control.	- Run the check of a single disorder;
	facts.		- Show the asserted facts;
	help.		- Show the list of available commands;
	clean.		- Clear the memory;
	quit.		- Close the interface.
\end{lstlisting}


\subsection{Clean}

The \texttt{clean} command clears the entire working memory. Unlike the implicit
clean-up performed at the beginning of every investigation or control --- which
removes only the proof-tree traces --- this command removes \emph{both} the proof
trees and all the facts asserted by the user during the session. It is therefore
used whenever the expert wants to start a brand-new consultation, with no
information carried over from previous runs:

\begin{lstlisting}
> clean.
\end{lstlisting}


\subsection{Quit}

The \texttt{quit} command closes the user interface and terminates the system.
Before halting, it also clears the working memory, so that no facts or proof
trees are left behind:

\begin{lstlisting}
> quit.
\end{lstlisting}
